% SYNTHETIC EVALUATION DRAFT — written for the retrieval eval in evals/.
% Near-topic distractor for diffusion.tex. Not a real paper; contains no results.
\documentclass{article}
\title{Stabilizing Large-Scale Generative Adversarial Networks for Text-Conditioned Image Synthesis}
\begin{document}
\maketitle

\begin{abstract}
Generative adversarial networks (GANs) synthesize an image in a single forward
pass of the generator, which makes them attractive for real-time and
interactive applications. Scaling GANs to large, diverse, text-conditioned
datasets has proven difficult because adversarial training becomes unstable
and suffers from mode collapse. This draft proposes a regularized training
recipe for a style-based generator with a projected multi-scale discriminator,
aiming to scale adversarial training to open-domain text-to-image generation.
\end{abstract}

\section{Introduction}
A GAN trains two networks in a minimax game: a generator maps a latent noise
vector to an image, and a discriminator tries to distinguish generated images
from real ones. At equilibrium the generator distribution matches the data
distribution. Style-based generators produce high-resolution, photorealistic
faces and objects on curated single-domain datasets, and sampling requires
only one network evaluation.

Training remains fragile. The discriminator can overpower the generator,
gradients vanish, and the generator may cover only a few modes of the data,
producing low diversity. These problems become worse as dataset diversity and
model size grow, and conditioning on free-form text adds a second objective:
generated images must not only look real but also match the caption.

Our contributions are: (1) a training recipe that combines gradient penalties
on real data, spectral normalization of the discriminator and adaptive
discriminator augmentation to keep the game balanced at scale; (2) a projected
discriminator that operates on features of a frozen pretrained vision encoder
at several resolutions, with a text-image matching head; and (3) a study of
how generator capacity and discriminator capacity should grow together.

\section{Related Work}
The original GAN formulation was followed by DCGAN, which introduced
convolutional architectures, and by Wasserstein GANs with gradient penalty,
which replaced the Jensen--Shannon objective with an Earth mover distance.
Spectral normalization constrains the discriminator's Lipschitz constant.
Progressive growing, BigGAN and the StyleGAN family improved fidelity through
architecture and scale; StyleGAN introduced a mapping network and per-layer
style modulation. Projected GANs use pretrained feature networks inside the
discriminator. Conditional GANs condition on class labels or text via
projection discriminators. Mode collapse has been studied through minibatch
discrimination and unrolled optimization. Likelihood-based and score-based
generative models are slower to sample and are not the focus here.

\section{Method}
\subsection{Generator}
The generator follows a style-based design: a mapping network transforms the
latent code and a pooled text embedding into a style vector, which modulates
the convolution weights at each resolution. Cross-attention layers at
intermediate resolutions attend to per-token text embeddings.

\subsection{Discriminator}
Real and generated images are encoded by a frozen pretrained vision backbone.
Lightweight discriminator heads operate on feature maps at four scales and
output a realism logit; a separate head scores image-text agreement.

\subsection{Regularization}
We apply an $R_1$ gradient penalty on real images, spectral normalization in
the discriminator heads, and adaptive augmentation whose probability is
adjusted from the discriminator's overfitting signal. The non-saturating
logistic loss is used for both networks.

\section{Experimental Plan}
We plan to train on a large image-caption dataset and evaluate at 64 and 256
pixel resolution. Metrics are Fréchet Inception Distance, precision and recall
for generative models to separate fidelity from diversity, CLIP score for text
alignment, and generation time per image. Baselines include StyleGAN-T,
GigaGAN-style models and projected GANs. Ablations remove each regularizer.

\section{Limitations}
Adversarial training requires careful hyperparameter tuning, and the frozen
feature network may bias the discriminator toward the features it was trained
on. Fine-grained text compositionality remains challenging for single-pass
generators.

\end{document}
